% Generated draft: requires mathematical and source review.
\documentclass[11pt]{article}
\usepackage[margin=25mm]{geometry}
\usepackage{fontspec}
\setmainfont{Latin Modern Roman}
\usepackage{amsmath,amssymb,mathtools,unicode-math}
\setmathfont{Latin Modern Math}
\usepackage{enumitem,needspace,xurl,hyperref}
\setlength{\parindent}{0pt}
\setlength{\parskip}{5pt}
\setlength{\emergencystretch}{3em}
\newcommand{\N}{\mathbb{N}}
\newcommand{\Z}{\mathbb{Z}}
\newcommand{\Q}{\mathbb{Q}}
\newcommand{\R}{\mathbb{R}}
\newcommand{\C}{\mathbb{C}}
\newcommand{\F}{\mathbb{F}}
\newcommand{\A}{\mathbb{A}}
\newcommand{\PP}{\mathbb{P}}
\newcommand{\E}{\mathbb{E}}
\newcommand{\eps}{\varepsilon}
\newcommand{\dd}{\,\mathrm{d}}
\begin{document}
\section*{UnsolvedMath ID: 10800005}
\textbf{Source problem code:} catalogue entry
\subsection{Definitions and mathematical statement}
Problem 1E asks for the obstructions to realizing the source's ``chains of formal changes'' by paths in the parameter space \(\R^k\). A path would be a map \(\gamma:[0,1]\to\R^k\), with the regularity, endpoints, allowed degeneracies and crossing rules prescribed by the underlying deformation problem. The record supplies no definition of those chains, no deformation family, no discriminant or stratification, and no rule assigning a chain to a path. Thus the phrase ``these chains'' has a missing antecedent essential to the mathematical target.

Once that data is fixed, the requested characterization is of which formal chains can arise from genuine parameter paths and what prevents realization of the others. It is not simply an assertion that \(\R^k\) is path connected: constraints from the unprovided singularity and allowed events determine realization.

The supplied research summary reports an upgraded enumeration algorithm applied to all simple and all parabolic singularities, with complete component classifications in those cases and no non-realizable formal chains found. It distinguishes those cases from arbitrary non-simple singularities, where it says a general no-obstruction theorem and the obstruction-characterization question remain open. These reports do not identify the underlying chain system or supply a proof of a universal no-obstruction claim; the literal Problem 1E remains the characterization request above.
\subsection{Short English statement}
Characterize when the source's formal change sequences can be realized by paths in its real parameter space, and identify the obstructions. The sequences and deformation rules are missing. The supplied progress assessment reports successful simple/parabolic classifications while leaving the general characterization unresolved.
\subsection{Sources}
\begin{itemize}
\item Local UnsolvedMath registry, permanent ID \(10800005\), exact source code \texttt{catalogue entry}, Problem 1E, attributed to Victor Vassiliev; no proposed year supplied. The title is truncated.
\item Vassiliev, \emph{A Few Problems on Monodromy and Discriminants}, Problem 1E: \url{https://amj.math.stonybrook.edu/html-articles/Files-2015-2024/15-11/}. The supplied background describes eighteen explicitly labeled theorem-style problem environments.
\item HTML-structure extraction, accessed 2026-07-29. Source evidence NEEDS\_REVIEW presents an open problem and calls for a current primary-literature audit. Registry status partially\_solved and classification PARTIAL-PROGRESS report algorithmic progress and simple/parabolic component classifications, with the general non-simple claim and obstruction characterization still open. No separate algorithm/classification references are supplied.
\item Dataset attribution: Ulam AI / UnsolvedMath Contributors, CC BY 4.0. Source rights NEEDS\_REVIEW: publicly accessible author/publisher copy, redistribution terms requiring review.
\end{itemize}
\end{document}
